\documentclass[11pt,a4paper]{article}
\usepackage{fontspec,amsmath,amssymb,amsthm,microtype}
\setmainfont{lmroman10-regular.otf}[BoldFont=lmroman10-bold.otf,ItalicFont=lmroman10-italic.otf,BoldItalicFont=lmroman10-bolditalic.otf]
\usepackage[margin=25mm]{geometry}
\usepackage[unicode,hidelinks]{hyperref}
\setlength{\emergencystretch}{3em}
\setlength{\parskip}{0.4em}
\setlength{\parindent}{0pt}
\allowdisplaybreaks[2]
\raggedbottom
\title{\texorpdfstring{Explicit accuracy budgets and integer iteration counts}{Explicit accuracy budgets and integer iteration counts}}
\author{Omer Deniz Akyildiz\\[0.4em]\normalsize Prepared with AI assistance}
\date{}
\begin{document}\maketitle
Draft mathematical exposition; these notes may contain mistakes.


The paper's step-size and iteration formulas are scalar consequences of its numerical error estimate. This section verifies that implication, retains a stability-compatible bias budget, and rounds the logarithmic threshold to a natural number. It handles zero initial error separately and clips negative logarithmic thresholds at zero. These accuracy corollaries have an explicit error inequality as input; they do not replace the outstanding unconditional stochastic theorem.



The error theorem becomes an executable prescription only after selecting a positive step and an integer count. The logarithmic threshold handles initialization, while an explicit allocation controls the persistent bias.


\section{\texorpdfstring{Replace geometric decay by its exponential envelope}{Replace geometric decay by its exponential envelope}}


For \(\alpha\le1\), the base \(1-\alpha\) is nonnegative. The elementary exponential inequality gives \(1-\alpha\le e^{-\alpha}\). Raising this comparison to a natural power preserves order, so \((1-\alpha)^n\le e^{-\alpha n}\). The argument uses the inequality \(1+s\le e^s\), monotonicity of nonnegative powers, and the exponential power identity. The lemma itself needs no positive lower bound on \(\alpha\); that requirement enters when dividing the logarithmic threshold.



For \(0\le\alpha\le1\), \((1-\alpha)^n\le e^{-\alpha n}\). The exponential form can be solved explicitly for the required iteration number. Its use must retain the admissible range of \(\alpha\), which comes from the sampler’s step restrictions.


\section{\texorpdfstring{Round the logarithmic threshold and include small initial errors}{Round the logarithmic threshold and include small initial errors}}


For initial coupling radius \(D\), define

\[n_0=\left\lceil\max\left\{0,\frac{\log(24D/\varepsilon)}\alpha\right\}\right\rceil.\]

This is a natural number even when the logarithm is negative. If \(\alpha>0\) and \(n\ge n_0\), then \(\alpha n\ge\log(24D/\varepsilon)\). The proof combines the ceiling's lower bound, the right-hand inequality for a maximum, and the cast of the natural comparison, before multiplying by the positive \(\alpha\). The explicit maximum and ceiling handle the initialization-dependent small-error regime and integer rounding.



The count is a natural-number ceiling of a nonnegative threshold. Taking the maximum with zero ensures that an already small initial error requires no negative number of steps. The ceiling adds fewer than one extra iteration beyond the real-valued requirement.


\section{\texorpdfstring{Verify the transient budget only when logarithmic cancellation is valid}{Verify the transient budget only when logarithmic cancellation is valid}}


For \(D>0\) and \(\varepsilon>0\), exponentiating the iteration inequality gives \(e^{-\alpha n}\le e^{-\log(24D/\varepsilon)}=\varepsilon/(24D)\). The geometric envelope and multiplication by \(D\) imply \((1-\alpha)^nD\le\varepsilon/24\). One proves positivity of the logarithm argument before using the exponential–logarithm identity, rewrites the negative exponential as an inverse, and cancels the nonzero factors. Zero initial error is handled by the next theorem without a logarithmic cancellation.



When the logarithm is positive, multiplying the chosen threshold by the contraction rate cancels its denominator. Exponentiating then reduces the initial error to the allocated fraction of accuracy. When the logarithm is nonpositive, the initial error already satisfies that allocation.


\section{\texorpdfstring{Combine the manuscript's 0.\allowbreak{}94 bias allocation with the transient}{Combine the manuscript's 0.94 bias allocation with the transient}}


Suppose an explicit estimate gives \(E_n\le\sqrt2(1-\alpha)^nD+b\), with \(0<\alpha\le1\), \(D\ge0\), and \(b\le0.94\varepsilon\). When \(D=0\), the bias already gives the conclusion and there is no iteration requirement. Otherwise use the rounded threshold. The exact rational inequality \(\sqrt2\le36/25\) bounds the remaining contribution by \(0.06\varepsilon\), so the total error is at most \(\varepsilon\).



One proves the square-root comparison by squaring, splits on zero versus positive initial radius, derives the logarithmic condition in the positive case, and combines the inequalities with nonnegative scalar multiplication. For KLMC, substitute \(\alpha=3mh/(4\gamma)\) and bias \(Mh\sqrt{2p}/m\), with a step size that is the minimum of the stability threshold and bias allowance. The same scalar argument applies to KIPLMC1 with its thermal radius and a separately budgeted concentration term.



The persistent bias and the decaying transient consume separate portions of the requested accuracy. The specific numerical allocation leaves enough slack for the \(\sqrt2\) positional prefactor and the rounded count. Adding the two verified contributions gives the final epsilon guarantee.


\end{document}
